% Shared V5 benchmark chapter, included by both solver user guides.
\section{Scheduler benchmark and direct v4 baseline}\label{sec:benchmark}

\subsection{Experimental design and metrics}
\code{run_scheduler_benchmark_v5.m} uses MEX so the comparison isolates
scheduler order from kernel implementation. It tests all 20 bundled cases
through case57 under three policies and three repeats: 180 measured runs.
A 23-worker pool remains open, each policy receives an untimed warm-up, policy
order rotates by repeat, and raw state is saved after each run for resumability.
Every numerical and timing value published for v5 comes solely from official
dataset \code{20260724_3x23}, produced on the designated main job machine under
MATLAB R2022a on Windows x64. Results from other machines are not merged into
the tables, plots, or release claims. Machine-readable provenance and scheduler
settings are recorded in \code{../BENCHMARK_METADATA_v5.json}; the portable raw
evidence is distributed as a separately attached GitHub Release archive.
The designated-machine run used the source-equivalent V5.5 development tree;
V5 differs only by the reproducibly rebuilt KLU binary. The metadata records
both KLU hashes, and the authors adopt the dataset as the official V5 result.

Recorded fields include final solutions, total traces, trace indices to 50\%,
90\%, and 99\%, wall time to 90\%, exhaustive wall time, worker trace-time
statistics, scheduler selection time, cursor checks, novelty draws, maximum
residual, and solution-set agreement. Trace-to-90\% measures anytime discovery;
exhaustive wall time includes the coverage tail needed to establish that no
eligible pair remains.

\subsection{Overall measured results}
\begin{center}\small
\begin{tabular}{@{}lrrrrr@{}}
\toprule
Policy & Wall total (s) & Summed t90 & Selection (s) & Selection share & t90 wins\\
\midrule
\code{scan} & 606.95 & 35,975 & 3.57 & 0.59\% & 6\\
\code{bandit} & 607.18 & 16,611 & 0.98 & 0.16\% & 7\\
\code{novelty} & 606.70 & 11,975 & 34.16 & 5.63\% & 18\\
\bottomrule
\end{tabular}
\end{center}
Exhaustive totals differed by less than 0.05\% from \code{scan}; the main effect
was \emph{when} solutions appeared. Relative to scan, bandit reduced aggregate
trace-to-90\% by 53.8\% and novelty by 66.7\%. Bandit is default because it
retained most of the anytime benefit with far lower selection overhead. For
case57, mean t90 was 19,481, 6,996, and 4,594 for scan, bandit, and novelty,
while exhaustive wall time remained approximately 333--337 seconds.

\IfFileExists{../scheduler_benchmark_v5_wall.png}{
\begin{figure}[htbp]\centering
\includegraphics[width=0.96\textwidth]{../scheduler_benchmark_v5_wall.png}
\caption{Three-repeat mean exhaustive MEX wall time; error bars show one sample
standard deviation.}
\end{figure}}{}
\IfFileExists{../scheduler_benchmark_v5_t90.png}{
\begin{figure}[htbp]\centering
\includegraphics[width=0.96\textwidth]{../scheduler_benchmark_v5_t90.png}
\caption{Mean completed traces required to reach 90\% of the final solution count.}
\end{figure}}{}
\IfFileExists{../scheduler_benchmark_v5_anytime.png}{
\begin{figure}[htbp]\centering
\includegraphics[width=0.96\textwidth]{../scheduler_benchmark_v5_anytime.png}
\caption{Anytime discovery for case57mod and case57; bands span the repeats.}
\end{figure}}{}
\IfFileExists{../scheduler_benchmark_v5_anytime_representative.png}{
\begin{figure}[htbp]\centering
\includegraphics[width=0.96\textwidth]{../scheduler_benchmark_v5_anytime_representative.png}
\caption{Anytime discovery for case14mod, case33bw, case39, and case30. Lines
are three-repeat means and bands span the repeat range. The figure is derived
from the stored per-trace histories; it does not require new solver runs.}
\end{figure}}{}
\IfFileExists{../scheduler_benchmark_v5_anytime_suite.png}{
\begin{figure}[htbp]\centering
\includegraphics[width=0.96\textwidth]{../scheduler_benchmark_v5_anytime_suite.png}
\caption{Normalized anytime discovery across all 20 benchmark cases. Cases and
repeats receive equal weight, with interquartile bands. Trace normalization
(left) isolates scheduling progress; elapsed-time normalization (right) also
reflects trace-duration and scheduler-overhead effects.}
\end{figure}}{}

The representative panels confirm that adaptive dispatch can improve early
coverage outside the two largest cases, but the gain is case-dependent. The
suite-level figure is a macro average: each case contributes equally, rather
than weighting results by solution count or allowing case57 to dominate.

\subsection{Large-case reference connectivity}
The benchmark retained final solution matrices and cumulative trace histories
for case30 and case57, but not the policy-specific \code{VBook} arrays required
to reconstruct exact scan/bandit/novelty edges. The diagrams below are therefore
single reference traversal graphs reconstructed offline from the complete
released-v4 MEX snapshots; no solver was rerun. They describe those recorded
continuation connections, not a three-policy v5 connectivity comparison.

The underlying v4 solution matrices were compared symmetrically and
antipode-aware with the saved v5 scan references. Distances were
$9.45\times10^{-12}$ for case30 and $2.00\times10^{-10}$ for case57, both far
below the $4\times10^{-7}$ deduplication tolerance.

\IfFileExists{../connectivity_case30_v4_reference.png}{
\begin{figure}[p]\centering
\includegraphics[width=0.96\textwidth]{../connectivity_case30_v4_reference.png}
\caption{case30 released-v4 reference connectivity: 472 solutions, 2,943
unique edges, and one component. Color is reference discovery order, node size
is degree, and line width reflects the number of shared traces.}
\end{figure}}{}
\IfFileExists{../connectivity_case57_v4_reference.png}{
\begin{figure}[p]\centering
\includegraphics[width=0.96\textwidth]{../connectivity_case57_v4_reference.png}
\caption{case57 released-v4 reference connectivity: 1,322 solutions, 8,301
unique edges, and one component. Color is reference discovery order, node size
is degree, and line width reflects the number of shared traces.}
\end{figure}}{}

\subsection{Solution-set validation and checkpoint test harness}
All 180 timing runs returned expected counts. Symmetric nearest-column
$L_\infty$ comparison checked bandit and novelty against scan for every case:
40/40 passed. Worst set distance was $2.6591\times10^{-9}$ versus the
$4\times10^{-7}$ tolerance; largest validation residual was
$5.6091\times10^{-8}$. The release also includes
\code{test_checkpoint_policy_compatibility_v5.m} for users to exercise
cross-policy resume on their own platform. No checkpoint result produced on a
non-designated machine is included as published numerical evidence.

\subsection{Direct comparison with released v4}
The retained 2026.07.15 v4 pass used the same 20 cases, MATLAB R2022a,
Windows x64, and 23 workers:
\begin{center}\small
\begin{tabular}{@{}lrrr@{}}
\toprule
Implementation & Aggregate wall (s) & Trace wall (s) & Solutions\\
\midrule
MEX v4 & 718.656 & 716.062 & 3256\\
Pure MATLAB v4 & 1272.966 & 1271.222 & 3256\\
V5 MEX \code{scan}, three-repeat mean & 606.947 & -- & 3256\\
\bottomrule
\end{tabular}
\end{center}
The v5 scan result is descriptively 15.5\% below historical MEX v4, but v4
is one retained pass while v5 is a current three-repeat experiment. This is
release context, not a controlled paired speedup claim. Full per-case results
are in \code{../SCHEDULER_BENCHMARK_v5.md} and the published CSV files.
